\chapter{Censo de la Industria de la Investigación de Mercados y Opinión Pública 2019} \label{censo2019}

\textbf{Nota:} El texto siguiente corresponde a la Edición XII (2019-2019) del Estudio Anual de la Industria de Investigación de Mercados y Opinión Pública en México\footnote{Fuente: Asociación Mexicana de Agencias de Inteligencia de Mercado y Opinión A.C. (“AMAI”), a través del Centro de Ciencia de Datos ITAM}.\newline


\noindent Este CENSO ha sido promovido por la AMAI con la finalidad de contar con una estimación lo más certera posible del volumen y características de la industria de investigación de mercados y opinión pública en México.
A partir de los datos de este censo, la AMAI elaborará un documento del estado de nuestro mercado, mismo que será difundido a los participantes del ejercicio, los medios de comunicación y los interesados en México y el extranjero.

Dada la naturaleza de la información que se pide, este cuestionario está dirigido al Presidente, Director General o socio principal de la empresa. Los datos recabados serán mantenidos en confidencialidad.
Sólo se divulgarán públicamente datos agregados y no información específica de una empresa en particular.

En el caso de la facturación, según la política adoptada desde el año 2002, para la difusión pública se hará un rankeo, sin especificar cifras, de aquellas empresas que expresamente lo autoricen.\newline

I. CONFIDENCIALIDAD DE LOS DATOS\newline

1.- ¿Autoriza que el nombre de su empresa aparezca en el rankeo público de compañías de investigación, ubicando su lugar a partir del dato de facturación proporcionado, pero sin revelarlo públicamente? (obligatoria)

SELECCIONE SOLAMENTE UNA OPCIÓN
\begin{itemize}
    \item[1)] (\rule{20mm}{0.1mm}) Sí
    \item[2)] (\rule{20mm}{0.1mm}) No
\end{itemize}

2.- Nombre oficial de la empresa (obligatoria)\newline

3.- Nombre del informante (obligatoria)\newline

II. TIPOS DE INVESTIGACIÓN REALIZADA\newline
4.- ¿En 2019 su empresa realizó proyectos bajo los siguientes tipos de investigación? (obligatoria)

\begin{itemize}
    \item[A)] Cualitativos a partir de interacciones con informantes: incluye sesiones de grupo, entrevistas a profundidad, etnografía aplicada, ya sea presencial o remota (por ejemplo: en línea), observaciones de conductas de consumo o procesos de venta (compras simuladas), etcétera
    \item[B)] Cuantitativos a partir de entrevistas con informantes: incluye todo tipo de proyectos hechos a partir de entrevistas estructuradas con grupos amplios de informantes, seleccionados aleatoriamente o bajo otro tipo de procedimiento. Las entrevistas pudieron haberse hecho cara a cara o en forma remota (telefónicamente, en línea, etcétera)
    \item[C)] Recolección y analítica de datos, tanto primarios como secundarios. Incluye auditorías de ventas y análisis de compraventas realizadas por bases de datos de clientes de una compañía; análisis de big data y otros relacionados, modelación de conductas y procesos de compra para productos de un cliente, etcétera
    \item[D)] Aplicaciones de biometría o neurociencia para determinar respuestas a estímulos diversos o para apoyar la realización de proyectos encaminados a entender o generar una determinada decisión humana
    \item[E)] Servicios de consultoría, asesoría o capacitación que se hayan facturado independientemente de proyectos de investigación de mercados y opinión pública
\end{itemize}

III. FACTURACIÓN EN 2019\newline

DESCONTAR FACTURACIÓN CRUZADA ENTRE FILIALES Y SUBCONTRATACIÓN A OTRAS EMPRESAS DE INVESTIGACION QUE NO SON MAQUILADORES. EN CASO DE CORPORATIVOS, CONSOLIDAR LA FACTURACION DE TODAS LAS UNIDADES DE NEGOCIO

5.- Por favor escriba con número el TOTAL de la facturación que tuvo su empresa en 2019. NOTA: Utilizar solamente dígitos, sin comas, espacios, ni signos de pesos. (obligatoria)
1) \rule{20mm}{0.1mm} pesos mexicanos

6.- Por favor escriba el porcentaje aproximado de la facturación total que tuvo su empresa en 2019. NOTA: El total de los porcentajes deberá sumar 100\%. (obligatoria)
VERIFIQUE QUE LA SUMA SEA 100\%

\begin{itemize}
    \item[1)] \rule{20mm}{0.1mm} \% Del total en Proyectos CUALITATIVOS
\item[2)] \rule{20mm}{0.1mm} \% Del total en Proyectos CUANTITATIVOS
\item[3)] \rule{20mm}{0.1mm} \% Del total en RECOLECCIÓN Y ANALÍTICA DE DATOS
\item[4)] \rule{20mm}{0.1mm} \% Del total en APLICACIONES DE BIOMETRÍA O NEUROCIENCIA
\item[5)] \rule{20mm}{0.1mm} \% Del total en servicios de CONSULTORÍA, asesoría o capacitación
\item[6)] \rule{20mm}{0.1mm} \% Otros
\end{itemize}

\rule{20mm}{0.1mm} \% TOTAL

7- De acuerdo a su mejor estimación, por favor indique qué porcentaje de su facturación anual total en 2019 correspondió a los siguientes rubros:\newline

1) DESTINATARIO

\rule{20mm}{0.1mm} \% del total

\begin{itemize}
    \item[a)] Para clientes o peticionarios establecidos en México \rule{20mm}{0.1mm}
\item[b)] Para clientes o peticionarios de países extranjeros \rule{20mm}{0.1mm}
\end{itemize}

2) SECTOR

\begin{itemize}
    \item[a)] Sector público / gobierno de México u otro país
\item[b)] Sector privado de México u otro país \rule{20mm}{0.1mm}
\item[c)] ONGs o instituciones similares de México u otro país \rule{20mm}{0.1mm}
\end{itemize} 

3) INDUSTRIA

\begin{itemize}
    \item[a)] Productos de consumo perecederos \rule{20mm}{0.1mm}
\item[b)] Productos de consumo no perecederos \rule{20mm}{0.1mm}
\item[c)] Comercio y menudeo \rule{20mm}{0.1mm}
\item[d)] Servicios financieros \rule{20mm}{0.1mm}
\item[e)] Automotriz \rule{20mm}{0.1mm}
\item[f)] Farmaceútica \rule{20mm}{0.1mm}
\item[g)] Telecomunicaciones y tecnología de comunicación \rule{20mm}{0.1mm}
\item[h)] Medios de comunicación y entretenimiento \rule{20mm}{0.1mm}
\item[i)] Agencias de publicidad \rule{20mm}{0.1mm}
\item[j)] Otros \rule{20mm}{0.1mm}
\end{itemize} 
IV. VOLUMEN DE INVESTIGACIÓN \newline

INVESTIGACIÓN CUALITATIVA HECHA EN 2019 \newline

8. Por favor indique su mejor estimación sobre el volumen de investigación cualitativa hecha por su empresa en 2019: \rule{20mm}{0.1mm} (NO INCLUIR EN ESTE RECUENTO PROYECTOS HECHOS PARA OTRA EMPRESA DE INVESTIGACION DE MEXICO, PERO INCLUIR LAS QUE SE HICIERON COMO MAQUILA PARA CLIENTES EXTRANJEROS).\newline

TIPO DE METODOLOGÍA \newline
UNIDADES REALIZADAS EN 2019 

\begin{itemize}
    \item Sesiones de grupo presenciales o remotas: \rule{20mm}{0.1mm}
\item Entrevistas a profundidad: \rule{20mm}{0.1mm}
\item Etnografías o aplicaciones de antropología: \rule{20mm}{0.1mm}
\end{itemize} 

9. Número de proyectos de investigación cualitativa realizados en 2019: \rule{20mm}{0.1mm} \newline 

INVESTIGACIÓN CUANTITATIVA HECHA EN 2019 \newline

10- Por favor indique su mejor estimación sobre el volumen de investigación cuantitativa hecha por su empresa en 2019: \rule{20mm}{0.1mm} (NO INCLUIR EN ESTE RECUENTO PROYECTOS HECHOS PARA OTRA EMPRESA DE INVESTIGACION DE MEXICO, PERO INCLUIR LAS QUE SE HICIERON COMO MAQUILA PARA CLIENTES EXTRANJEROS)\newline

TIPO DE LEVANTAMIENTO\newline
NÚMERO DE UNIDADES REALIZADAS \newline

\begin{itemize}
    \item Cara a cara en domicilios (viviendas o empresas) \rule{20mm}{0.1mm}
    \item Telefónico \rule{20mm}{0.1mm}
    \item En vía pública / punto de afluencia \rule{20mm}{0.1mm}
    \item A la salida de casilla electoral \rule{20mm}{0.1mm}
    \item Por internet (con personas o empresas) \rule{20mm}{0.1mm}
    \item Otro \rule{20mm}{0.1mm}
\end{itemize}

11.- Número de proyectos cuantitativos que su empresa realizó en 2019, (obligatoria)
      1) \rule{20mm}{0.1mm}  Proyectos cuantitativos con entrevistas\newline
      
V. RECURSOS HUMANOS Y TÉCNICOS \newline

12A.- ¿Cuántas personas colaboraron en su empresa en 2019 de forma permanente, POR NOMINA dentro de los siguientes rubros? (obligatoria) \rule{20mm}{0.1mm} \newline

12B.- De ese personal de planta, que se le paga POR NOMINA, ¿qué porcentajes corresponden a mujeres y a hombres? (haga su mejor estimación)\newline

CATEGORÍA

\begin{itemize}
    \item[1)] Socios y altos ejecutivos (directores generales y directores de área) NÚMERO TOTAL \rule{20mm}{0.1mm} \% MUJERES, \rule{20mm}{0.1mm} \% HOMBRES
    \item [2)] Personal técnico y de atención a clientes (subdirectores, investigadores, analistas, etc.) NÚMERO TOTAL \rule{20mm}{0.1mm} \% MUJERES, \rule{20mm}{0.1mm} \% HOMBRES
    \item [3)] Personal de apoyo administrativo (contadores, secretarias, mensajeros, choferes, etc.) NÚMERO TOTAL \rule{20mm}{0.1mm} \% MUJERES, \rule{20mm}{0.1mm} \% HOMBRES
    \item [4)] Equipo de producción (entrevistadores, capturistas, codificadores) NÚMERO TOTAL \rule{20mm}{0.1mm} \% MUJERES, \rule{20mm}{0.1mm} \% HOMBRES
\end{itemize}
 
GRAN TOTAL \rule{20mm}{0.1mm}


13A.- ¿Cuántas personas colaboraron en su empresa en 2019 eventualmente, POR PROYECTO dentro de los siguientes rubros? (obligatoria) \rule{20mm}{0.1mm}  \newline

13B.- De ese personal eventual, que se le paga POR PROYECTO, ¿qué porcentajes corresponden a mujeres y a hombres? (haga su mejor estimación)

CATEGORÍA

\begin{itemize}
    \item[1)] Socios y altos ejecutivos (directores generales y directores de área) NÚMERO TOTAL \rule{20mm}{0.1mm} \% MUJERES, \rule{20mm}{0.1mm} \% HOMBRES
    \item [2)] Personal técnico y de atención a clientes (subdirectores, investigadores, analistas, etc.) NÚMERO TOTAL \rule{20mm}{0.1mm} \% MUJERES, \rule{20mm}{0.1mm} \% HOMBRES
    \item [3)] Personal de apoyo administrativo (contadores, secretarias, mensajeros, choferes, etc.) NÚMERO TOTAL \rule{20mm}{0.1mm} \% MUJERES, \rule{20mm}{0.1mm} \% HOMBRES
    \item [4)] Equipo de producción (entrevistadores, capturistas, codificadores) NÚMERO TOTAL \rule{20mm}{0.1mm} \% MUJERES, \rule{20mm}{0.1mm} \% HOMBRES
\end{itemize}
 
GRAN TOTAL \rule{20mm}{0.1mm}

VI. EXPECTATIVAS Y PROSPECTIVA SOBRE LA INDUSTRIA

14.-Señale los cambios más importantes que vislumbra para nuestra industria de aquí a 5 años.

\begin{itemize}
    \item[1)] \rule{20mm}{0.1mm}
    \item [2)] \rule{20mm}{0.1mm}
    \item [3)] \rule{20mm}{0.1mm}
\end{itemize}

15.-¿Y cuáles piensa serán los retos más relevantes a confrontar por la industria en los próximo 5 años?
\begin{itemize}
    \item[1)] \rule{20mm}{0.1mm}
    \item [2)] \rule{20mm}{0.1mm}
    \item [3)] \rule{20mm}{0.1mm}
\end{itemize}

16.- En términos generales: ¿qué pronóstico tiene respecto a este año en comparación con el anterior para...?

2020 será... 
Mejor que 2019 Igual que 2019 Peor que 2019
\begin{itemize}
    \item Mi economía personal: \rule{20mm}{0.1mm}
    \item Mi empresa: \rule{20mm}{0.1mm}
    \item La industria en México: \rule{20mm}{0.1mm}
    \item La industria en el mundo: \rule{20mm}{0.1mm}
    \item México en general: \rule{20mm}{0.1mm}
\end{itemize}

17- Por último: AMAI está diseñando próximas actividades relacionadas con la innovación y la capacitación. Al respecto:
\begin{itemize}
    \item ¿Qué temas de innovación en la industria sugiere ser tratados prioritariamente en la agenda de AMAI? \rule{20mm}{0.1mm}
\item ¿Qué temas deberían incorporarse a los cursos y sesiones de capacitación de AMAI (por ejemplo, en Campus AMAI)? \rule{20mm}{0.1mm}

\end{itemize}

18.- Muchas gracias por su participación. Si desea agregar un comentario adicional escríbalo a continuación: \rule{20mm}{0.1mm}